\documentclass[10pt]{article}
\usepackage[letterpaper,margin=0.8in]{geometry}
\usepackage{amsmath}
\usepackage[T1]{fontenc}
\usepackage{lmodern}
\usepackage{xcolor}
\definecolor{ink}{HTML}{343B40}
\color{ink}
\setlength{\parindent}{0pt}
\setlength{\parskip}{6pt}
\title{Separating seasonal structure\\from gradual change}
\author{Atelier demonstration manuscript}
\date{Synthetic research note}
\begin{document}
\maketitle
\begin{abstract}
A time series can combine a gradual trend, periodic structure, and short-term variation. This fictional research note illustrates how these components can be represented in one linear model. Using 240 generated observations, we describe the model, quantify uncertainty under a known noise process, and distinguish in-sample agreement from predictive evaluation. All observations and results are synthetic.
\end{abstract}
\section{Research question}
How can a seasonal model distinguish gradual change from a repeating cycle? The example uses a regular observation schedule over two years. Its purpose is to make the relationship between a research question, an analysis script, and a scientific manuscript explicit. It does not describe a real population or study site.
\section{Model and assumptions}
Let $t_i$ denote time in months and $y_i$ a response in arbitrary units. We specify
\begin{equation}
y_i = \beta_0 + \beta_1\frac{t_i}{24}
+ \beta_2\sin\left(\frac{2\pi t_i}{12}\right)
+ \beta_3\cos\left(\frac{2\pi t_i}{12}\right) + \epsilon_i.
\end{equation}
The errors are independent Gaussian draws with a standard deviation of $0.12$, known by construction. The generating coefficients are $(0.35,\,0.28,\,0.42,\,-0.16)$. A fixed random seed makes the demonstration reproducible. These assumptions define an illustration, not a model validated for an empirical application.
\section{Estimation and uncertainty}
We estimate the coefficients by ordinary least squares. With design matrix $X$ and known noise variance $\sigma^2$, the sampling covariance is
\begin{equation}
\operatorname{Var}(\hat{\boldsymbol\beta})=\sigma^2(X^\mathsf{T}X)^{-1}.
\end{equation}
At an observation coordinate with design vector $\mathbf{x}_*$, uncertainty in the fitted mean follows from $\mathbf{x}_*^\mathsf{T}\operatorname{Var}(\hat{\boldsymbol\beta})\mathbf{x}_*$. The plotted intervals are pointwise 95\% confidence intervals for the mean. They are not prediction intervals for individual observations.
\section{Interpretation and limits}
The accompanying figure brings together the observations, fitted mean, estimated coefficients, and residuals. Agreement between observed and fitted values is evaluated on the same data used to estimate the model. It should therefore not be interpreted as evidence of performance on new observations.

A real study would require justified sampling assumptions, appropriate references, and validation suited to its scientific question. Here, the manuscript and its figures serve only to demonstrate reading, annotation, analysis, and revision within Atelier.
\section*{Demonstration materials}
\small
[1] \textit{Synthetic seasonal regression}. Companion script \texttt{analysis.py}, seed 2026.\par
[2] \textit{Seasonal dynamics, resolved}. Companion four-panel figure generated from the same model.\par
\textcolor{gray}{All materials are fictional. No empirical findings or external publications are claimed.}
\end{document}
